% Chapter 8: Orbit Closure and the Active Adversary
% DEFINES: ch:orbit-closure; sec:active-adversary, sec:orbit-closure-def,
%   sec:residual-entropy, sec:typed-chains, sec:ch8-notes;
%   def:orbit, thm:monotonicity, thm:orbit-bound, cor:orbit-set-form,
%   ex:orbit-names, rem:active-vs-marginal
% REFERENCES (backward): def:uniformity, ch:boolean-algebra, thm:not-approximate,
%   sec:granularity
% REFERENCES (forward, part labels): part:confidentiality

\chapter{Orbit Closure and the Active Adversary}
\label{ch:orbit-closure}

Every confidentiality statement so far has been about a \emph{passive}
observer: the untrusted machine holds tokens and studies their distribution,
and $\delta$ bounds what it can read from a single one. But the untrusted
machine in this framework does more than watch. It \emph{evaluates}: it
holds cipher maps, the Boolean operations of \Cref{ch:boolean-algebra} among
them, and can apply them to the tokens it has, building new tokens and
watching what coincides. This chapter measures what such an active adversary
can learn, and the measure is a new one. It is not $\delta$, and the
distinction between the two is the boundary between this part and the
confidentiality theory that follows.

\section{The Active Adversary}
\label{sec:active-adversary}

The shift is from observation to operation. A passive adversary sees a token
$c$ and asks what its distribution leaks; the answer is governed by
representation uniformity, $\delta$. An active adversary holds a set of
cipher maps and applies them to $c$ in every combination, accumulating the
tokens it can reach and noting which computations land on the same value.
What it learns is no longer a property of $c$'s marginal distribution; it is
a property of how far the exposed operations let it travel from $c$. A token
can be marginally near-uniform, $\delta \approx 0$, and still surrender its
latent value to an adversary whose operation set is rich enough to separate
it from its neighbors. Representation uniformity does not bound this, and the
quantity that does is combinatorial: the size of the reachable set.

\section{Orbit Closure}
\label{sec:orbit-closure-def}

\begin{definition}[Orbit closure]
\label{def:orbit}
Let $F = \{\fhat_1, \ldots, \fhat_k\}$ be the cipher maps available to the
untrusted machine, and $c \in \B^n$ a token. The \emph{orbit closure}
$\orbitF(c) \subseteq \B^n$ is the smallest set containing $c$ and closed
under the maps in $F$: if $c' \in \orbitF(c)$ and $\fhat_i \in F$ then
$\fhat_i(c') \in \orbitF(c)$. Equivalently, it is the set of all tokens
reachable from $c$ by composing operations of $F$ in any order.
\end{definition}

Because $\B^n$ is finite the orbit is finite, computed by breadth-first
expansion from $c$. Two features make it the right object. First, the orbit
\emph{set} depends only on $F$ and $c$, not on the latent function, the
domain, or the trapdoor; the designer who knows which operations to expose
can compute it without any secret, and read off in advance how far the
untrusted machine can roam. Second, it is exactly the support of what the
adversary observes: probing $c$ with $F$ resolves the computation to one of
the tokens in $\orbitF(c)$, and what carries information about the latent
value is \emph{which} one.

\begin{theorem}[Monotonicity]
\label{thm:monotonicity}
If $F \subseteq G$ then $\orbitF(c) \subseteq \mathrm{orbit}_G(c)$ for every
$c$.
\end{theorem}

The proof is a one-line induction on the orbit's construction: every token
reachable using $F$ is reachable using the larger $G$, since each step
available under $F$ is available under $G$. The reading is the design rule:
exposing more operations can only enlarge the reachable set, never shrink
it; with no operations at all, $\mathrm{orbit}_{\varnothing}(c) = \{c\}$ and
the adversary learns nothing beyond the bare token. Every operation the
trusted machine grants is an edge added to the reachability graph the
adversary explores.

\section{The Residual-Entropy Bound}
\label{sec:residual-entropy}

The orbit size turns directly into a confidentiality bound. Let $X$ be the
latent value encoded by $c$, drawn from the input distribution, and let the
adversary's observation, which token in $\orbitF(c)$ the computation
resolves to, be a random variable supported on that orbit.

\begin{theorem}[Confidentiality bound]
\label{thm:orbit-bound}
With $X$ the latent value behind $c$ and the active adversary's observation
supported on $\orbitF(c)$,
\[
  H(X \mid \text{view}) \;\geq\; H(X) - \log_2 \lvert \orbitF(c)\rvert .
\]
\end{theorem}

\begin{proof}
The observation is supported on $\orbitF(c)$, a set of size
$\lvert\orbitF(c)\rvert$, so its entropy is at most
$\log_2\lvert\orbitF(c)\rvert$. Mutual information is bounded by the entropy
of either side, so the information the view carries about $X$ is
$I(X; \text{view}) \le \log_2\lvert\orbitF(c)\rvert$. Substituting into
$H(X \mid \text{view}) = H(X) - I(X; \text{view})$ gives the claim.
\end{proof}

The bound has an equivalent set form, and the set form is where a recurring
error must be avoided. Writing the confidentiality as the fraction of the
latent space the adversary cannot rule out,
\[
  \mathrm{conf} \;\geq\; 1 - \frac{\lvert\orbitF(c)\rvert}{\lvert X\rvert},
\]
the denominator is $\lvert X\rvert$, the number of latent values, \emph{not}
$2^n$, the size of the token space. The orbit lives in $\B^n$, but it is
measured against the domain it could be encoding: an orbit of $m$ tokens can
narrow $X$ to at most $m$ possibilities out of $\lvert X\rvert$, so the bound
saturates, becomes vacuous, exactly when $\lvert\orbitF(c)\rvert \ge \lvert
X\rvert$, a rich enough operation set identifying every latent value.
Measuring the orbit against $2^n$ instead would make the bound trivially
satisfied for every realizable orbit and say nothing.

\begin{example}[Naming the four]
\label{ex:orbit-names}
Take the four names of \Cref{ex:watchlist-numbers} as the latent space,
$\lvert X\rvert = 4$ and $H(X) = 2$ bits under a uniform prior, and a token
$c$ encoding one unknown name. If the exposed operations generate an orbit
of size $\lvert\orbitF(c)\rvert = 2$, the bound gives a residual
$H(X \mid \text{view}) \ge 2 - \log_2 2 = 1$ bit, and the set form
$\mathrm{conf} \ge 1 - 2/4 = 0.5$: the adversary can halve the candidate
names but no more. Enrich the operations until the orbit reaches all four
tokens for the four names, and the bound goes to $2 - \log_2 4 = 0$: vacuous,
the name is in principle pinned down. The denominator throughout is the four
names, not the $4096$ strings of the token space; the orbit is read against
what it could mean, not against where it lives.
\end{example}

\begin{remark}[Which scale this is]
\label{rem:active-vs-marginal}
This bound is the \emph{active} scale, and it is governed by orbit size, not
by $\delta$. The two scales are independent in both directions: a token can
be marginally near-uniform ($\delta \approx 0$) and still leak its value
through a large orbit when the operation set is rich, and a small orbit can
cap active leakage even when $\delta$ is not negligible. No value of $\delta$
bounds the orbit, and no orbit bound improves $\delta$; they are different
axes and must not be conflated. The bound is also loose, deliberately: it
counts the orbit's size but ignores the functional relations among its
elements, which constrain the latent value further. Tightening it for
specific operation sets, and the matching lower bounds for composition, is
the work of \Cref{part:confidentiality}.
\end{remark}

\section{Typed Composition Chains}
\label{sec:typed-chains}

Orbit closure is the active face of a single cipher value; the compositional
scale appears when typed cipher maps are chained. Feeding the output of one
typed map into the next, the reachable set grows along the chain: a step
that combines $a$ prior values into one expands the frontier by the step's
arity, and the count of reachable tokens compounds through the chain rather
than merely adding. The total an adversary can reach is bounded by the sum
of these per-step frontiers, a quantity that grows with both the chain's
length and the arity of its operations.

What the orbit bound previews here, the full confidentiality theory
delivers. When two cipher maps share an input, or a chain exposes the joint
behavior of several latent values, the adversary's task becomes estimating a
\emph{joint} distribution, and the difficulty of that estimation, the number
of observations it takes and the precision it can reach, is the compositional
confidentiality scale. Its headline is a recovery rate of order
$\lvert Y_1\rvert\,\lvert Y_2\rvert / \xi^2$ in two latent variables, with a
matching lower bound that no per-map parameter, $\delta$ included, can move.
That rate, its tight lower bound, and the active-adversary measures that
accompany it are \Cref{part:confidentiality}; this chapter has shown why
they are needed, that operating on cipher values is a leakage channel of its
own, distinct from observing them, and not closed by representation
uniformity.

\section{Notes and Provenance}
\label{sec:ch8-notes}

Orbit closure, its monotonicity, and the residual-entropy bound are the
algebraic-cipher-types paper's, restated here with the elementary bound
proved in full and monotonicity sketched, per the part's proof-depth
decision. The set form $\mathrm{conf} \ge 1 - \lvert\orbitF(c)\rvert /
\lvert X\rvert$ is stated with the denominator $\lvert X\rvert$, the latent
space, guarding against the recurring slip of normalizing by the token space
$2^n$. The two-scale separation, active orbit leakage versus marginal
$\delta$, is the spine's, and it is the hinge between this part and the next:
\Cref{part:confidentiality} takes the marginal scale to its tight
entropy-ratio bound and the active scale to its compositional lower bounds,
and the warning of \Cref{rem:active-vs-marginal}, that no single parameter
governs both, is the reason that part needs the machinery it does. With
orbit closure the algebra of Part~III is complete: the operations
(\Cref{ch:boolean-algebra}), the types (\Cref{ch:algebraic-types}), and now
the leakage that operating on cipher values incurs.
